\section*{Итоги решения поставленных задач}

Цель работы, заключавшаяся в разработке системы управления контентом для корпоративного
маркетингового сайта на базе \texttt{Strapi} и \texttt{Astro}, достигнута в объеме,
зафиксированном для данной ВКР. В ходе исследования была проанализирована предметная
область маркетинговых сайтов, показаны ограничения подхода, при котором контент и
публичная витрина остаются жестко связаны, и обоснован выбор \texttt{CMS-first}
архитектуры с разделением контентного и frontend-слоев. На этой основе были сформированы
функциональные и нефункциональные требования к системе, а также спроектированы ее
архитектура, модель данных и редакторские сценарии.

В практической части работы реализован контур управления ключевыми сущностями через
\texttt{Strapi}, включая \texttt{home-page}, \texttt{page}, \texttt{article},
\texttt{project}, \texttt{vacancy}, связанные справочники и прикладные сущности форм.
Для маркетинговых страниц разработан конструктор на основе \texttt{Dynamic Zone}, а
frontend-витрина на \texttt{Astro} доведена до \texttt{CMS-first} модели с
локализованными маршрутами, \texttt{preview mode}, \texttt{SEO}/\texttt{Open Graph},
\texttt{sitemap} и publication-связкой \texttt{webhook -> rebuild}. Тем самым поставленные
задачи по анализу, проектированию, реализации и проверке решения были последовательно
решены и связаны в единый инженерный результат.

\section*{Практический результат работы}

Практическим результатом ВКР стала рабочая \texttt{CMS-first} платформа, в которой
\texttt{Strapi} выступает ядром управления контентом, а \texttt{Astro} --- публичной
витриной корпоративного маркетингового сайта. Контент-менеджер или маркетолог получает
возможность управлять главной страницей, маркетинговыми \texttt{slug}-страницами,
статьями, проектами, \texttt{SEO}-метаданными и глобальными данными сайта через
административный интерфейс, не изменяя исходный код frontend-приложения. Для карьерного
контура реализованы отдельные сущности вакансий, справочников и откликов, а прикладные
формы вынесены в server-side маршруты frontend с повторной валидацией и защитой
публичного ввода.

Существенным достоинством полученного результата является то, что он подтвержден не только
структурой кода, но и воспроизводимой доказательной базой. На актуальном baseline
зафиксированы успешная сборка frontend и \texttt{CMS}, генерация \texttt{37}
предсобранных публичных маршрутов, наличие \texttt{sitemap}, корректная работа
\texttt{preview}-сценариев, \texttt{SEO}-контура, локализованных маршрутов
\texttt{ru/en}, публичных форм и \texttt{webhook}-связки публикации. Автоматизированная
smoke-проверка подтвердила обязательные сценарии без отказов, а браузерный аудит показал
стабильное выполнение representative маршрутов без ошибок исполнения. Это позволяет
рассматривать результат работы как завершенное прикладное решение для рассматриваемой
предметной области, а не как набор разрозненных технических экспериментов.

\section*{Ограничения текущей версии}

Несмотря на достигнутую целостность решения, итоговая версия системы осознанно ограничена
рамками ВКР. Публичный двуязычный production-контур охватывает storefront-core,
маркетинговые страницы, статьи и проекты, тогда как карьерный модуль сохраняет отдельную
маршрутную границу и не переводится полностью в locale-prefixed публичную схему. Для
разделов-списков \texttt{articles}, \texttt{projects} и \texttt{vacancies} используется
route-owned слой метаданных, тогда как editor-managed \texttt{SEO} в обязательном объеме
доведено до \texttt{home-page}, \texttt{page} и detail-страниц сущностей
\texttt{article/project/vacancy}.

Дополнительно следует учитывать, что публикационный контур доказан локально как
versioned, config-validated и привязанный к managed \texttt{webhook}, однако внешний
\texttt{Dokploy}-segment не воспроизводился в рамках работы как полный live
\texttt{rebuild/redeploy}. Тестовый контур при этом построен как приемочный и
интеграционный baseline, а не как полный набор unit- и end-to-end тестов для всех
возможных ветвей поведения. Эти ограничения не обесценивают результат ВКР, но задают
корректную границу его доказанности и предотвращают завышение фактически подтвержденных
свойств системы.

\section*{Направления дальнейшего развития}

Логичным продолжением работы является расширение набора \texttt{Dynamic Zone}-блоков,
углубление редакторских workflow и развитие эксплуатационной автоматизации. В прикладном
плане это включает расширение мультиязычного публичного контура карьерного модуля,
дальнейшее развитие ролей и редакторских сценариев, а также подключение внешних
интеграций для обработки лидов и откликов. Отдельным направлением является усиление
production-доказательства deployment-контура за пределами локального baseline, включая
внешний \texttt{rebuild/redeploy} и дополнительные эксплуатационные метрики.

Таким образом, выполненная работа создает завершенную архитектурную и программную основу
для дальнейшего развития корпоративной маркетинговой платформы. Главный итог ВКР состоит
в том, что управление контентом, публикация, локализация, \texttt{preview},
\texttt{SEO}-контур и базовая эксплуатационная модель были сведены в единую
\texttt{CMS-first} систему, пригодную для практического использования и дальнейшего
масштабирования.
